\documentclass[../../../../main.tex]{subfiles}
\begin{document}

$a$は正の実数とする．
複素数$z$が$|z-1|=a$かつ$\displaystyle z\neq\frac{1}{2}$を満たしながら動くとき，
複素数平面上の点$\displaystyle w=\frac{z-3}{1-2z}$が描く図形をKとする．
このとき，次の問いに答えよ．
\begin{enumerate}
  \itemKが円となるための$a$の条件を求めよ．
また，そのときKの中心が表す複素数とKの半径を，それぞれ$a$を用いて表せ．
  \item$a$が(1)の条件を満たしながら動くとき，
虚軸に平行で円Kの直径となる線分が通過する領域を複素数平面上に図示せよ．
\end{enumerate}

\end{document}
